\documentclass[12pt]{article}

\usepackage{graphicx}
\usepackage{amsmath}
\usepackage{natbib}
\usepackage{hyperref}

\title{Effects of Noise Pollution on Marine Mammals: Physiological Stress Responses, Behavioral Changes, and Conservation Implications}
\author{Your Name \\ Department of Marine Biology, Your University}

\begin{document}

\maketitle

\begin{abstract}
This paper investigates the effects of noise pollution on marine mammals, focusing on physiological stress responses, behavioral changes, and the broader implications for conservation. Field studies and experimental data are reviewed to ensure a comprehensive understanding of these impacts. Conservation strategies and policy recommendations are discussed to propose a path forward for mitigating the negative consequences of marine noise pollution.
\end{abstract}

\section{Introduction}
The increasing levels of noise pollution in marine environments, primarily due to shipping, industrial activities, and naval exercises, pose significant threats to marine mammals. These mammals rely heavily on sound for communication, navigation, and finding food. Disruptions in their acoustic environment can lead to profound physiological and behavioral changes, ultimately affecting their survival and reproduction. This paper aims to explore these impacts and suggest effective conservation strategies.

\section{Impact Analysis}
\subsection{Physiological Stress Responses}
Noise pollution can induce a range of physiological stress responses in marine mammals. Elevated cortisol levels, increased heart rate, and other stress markers have been documented in several species. More detailed analyses from field studies should be incorporated here \citep{SampleFieldStudy2020}.

\subsection{Behavioral Changes}
Behavioral alterations include changes in vocalization patterns, avoidance of noisy areas, and disrupted feeding and mating behaviors. Long-term exposure may result in habitat displacement and reduced reproductive success \citep{ExperimentalStudy2021}.

\section{Case Studies}
\subsection{Case Study 1: Species A}
Provide details on specific studies or experimental data related to Species A, showcasing how noise pollution has impacted their physiology and behavior \citep{FieldStudySpeciesA2019}.

\subsection{Case Study 2: Species B}
Discuss another species to illustrate different aspects of the problem through a comparative analysis of different environments or noise sources \citep{FieldStudySpeciesB2020}.

\section{Conservation Strategies}
Effective conservation strategies must address the sources of noise pollution and its management. This involves international cooperation, strict regulations on noise levels, and the implementation of quieter technologies.
\subsection{Regulatory Aspects}
Discuss current regulations and policies, their effectiveness, and potential areas for improvement.

\subsection{Technological Approaches}
Highlight advancements in technology that can reduce noise pollution, such as quieter ship designs and better acoustic mitigation techniques.

\section{Recommendations}
Provide actionable recommendations based on the analysis. These should be focused on policy changes, further research needs, and practical implementations for reducing noise pollution and protecting marine mammals.

\section{Conclusion}
Summarize the key findings and emphasize the importance of taking immediate and effective actions to mitigate noise pollution in marine environments to ensure the health and sustainability of marine mammal populations.

\bibliographystyle{apalike}
\bibliography{prompt10_4o}

\end{document}